\documentclass[10pt,a4paper]{amsart}
\usepackage{amsmath,amssymb,amsthm,graphicx}
\usepackage{verbatim}
\newenvironment{reference}{\comment}{\endcomment}
%\newenvironment{reference}{}{}
\newenvironment{slogan}{\comment}{\endcomment}
\newenvironment{history}{\comment}{\endcomment}
% For commutative diagrams we use Xy-pic
\usepackage[all]{xy}

% We use 2cell for 2-commutative diagrams.
\xyoption{2cell}
% Theorem environments.
%
\theoremstyle{plain}
\newtheorem{theorem}[subsection]{Theorem}
\newtheorem{proposition}[subsection]{Proposition}
\newtheorem{lemma}[subsection]{Lemma}

\theoremstyle{definition}
\newtheorem{definition}[subsection]{Definition}
\newtheorem{example}[subsection]{Example}
\newtheorem{exercise}[subsection]{Exercise}
\newtheorem{situation}[subsection]{Situation}

\theoremstyle{remark}
\newtheorem{remark}[subsection]{Remark}
\newtheorem{remarks}[subsection]{Remarks}

\numberwithin{equation}{subsection}

% Macros
%
\def\lim{\mathop{\mathrm{lim}}\nolimits}
\def\colim{\mathop{\mathrm{colim}}\nolimits}
\def\Spec{\mathop{\mathrm{Spec}}}
\def\Hom{\mathop{\mathrm{Hom}}\nolimits}
\def\Ext{\mathop{\mathrm{Ext}}\nolimits}
\def\SheafHom{\mathop{\mathcal{H}\!\mathit{om}}\nolimits}
\def\SheafExt{\mathop{\mathcal{E}\!\mathit{xt}}\nolimits}
\def\Sch{\mathit{Sch}}
\def\Mor{\mathop{\mathrm{Mor}}\nolimits}
\def\Ob{\mathop{\mathrm{Ob}}\nolimits}
\def\Sh{\mathop{\mathit{Sh}}\nolimits}
\def\NL{\mathop{N\!L}\nolimits}
\def\CH{\mathop{\mathrm{CH}}\nolimits}
\def\proetale{{pro\text{-}\acute{e}tale}}
\def\etale{{\acute{e}tale}}
\def\QCoh{\mathit{QCoh}}
\def\Ker{\mathop{\mathrm{Ker}}}
\def\Im{\mathop{\mathrm{Im}}}
\def\Coker{\mathop{\mathrm{Coker}}}
\def\Coim{\mathop{\mathrm{Coim}}}

% Boxtimes
%
\DeclareMathSymbol{\boxtimes}{\mathbin}{AMSa}{"02}

%
% Macros for moduli stacks/spaces
%
\def\QCohstack{\mathcal{QC}\!\mathit{oh}}
\def\Cohstack{\mathcal{C}\!\mathit{oh}}
\def\Spacesstack{\mathcal{S}\!\mathit{paces}}
\def\Quotfunctor{\mathrm{Quot}}
\def\Hilbfunctor{\mathrm{Hilb}}
\def\Curvesstack{\mathcal{C}\!\mathit{urves}}
\def\Polarizedstack{\mathcal{P}\!\mathit{olarized}}
\def\Complexesstack{\mathcal{C}\!\mathit{omplexes}}
% \Pic is the operator that assigns to X its picard group, usage \Pic(X)
% \Picardstack_{X/B} denotes the Picard stack of X over B
% \Picardfunctor_{X/B} denotes the Picard functor of X over B
\def\Pic{\mathop{\mathrm{Pic}}\nolimits}
\def\Picardstack{\mathcal{P}\!\mathit{ic}}
\def\Picardfunctor{\mathrm{Pic}}
\def\Deformationcategory{\mathcal{D}\!\mathit{ef}}
\usepackage{cleveref}
\crefname{originaldefinition}{original definition}{original definitions}
\Crefname{originaldefinition}{Original definition}{Original definitions}
\makeatletter
\newlabel{definition-defined-on-perfect}{{0F9E}{source}{}{sourcepointer-1}{}}
\newlabel{definition-defined-on-perfect@cref}{{[originaldefinition][0][]0F9E}{[1][0][]source}{}{}{}}
\makeatother
\crefname{originallemma}{original lemma}{original lemmas}
\Crefname{originallemma}{Original lemma}{Original lemmas}
\makeatletter
\newlabel{lemma-base-change-envelope}{{0GU7}{source}{}{sourcepointer-2}{}}
\newlabel{lemma-base-change-envelope@cref}{{[originallemma][0][]0GU7}{[1][0][]source}{}{}{}}
\makeatother
\crefname{originalproposition}{original proposition}{original propositions}
\Crefname{originalproposition}{Original proposition}{Original propositions}
\makeatletter
\newlabel{more-morphisms-proposition-envelope-with-resolution-property}{{0GTU}{source}{}{sourcepointer-3}{}}
\newlabel{more-morphisms-proposition-envelope-with-resolution-property@cref}{{[originalproposition][0][]0GTU}{[1][0][]source}{}{}{}}
\makeatother
\crefname{originallemma}{original lemma}{original lemmas}
\Crefname{originallemma}{Original lemma}{Original lemmas}
\makeatletter
\newlabel{perfect-lemma-resolution-property-ample-family}{{0GMM}{source}{}{sourcepointer-4}{}}
\newlabel{perfect-lemma-resolution-property-ample-family@cref}{{[originallemma][0][]0GMM}{[1][0][]source}{}{}{}}
\makeatother
\crefname{originallemma}{original lemma}{original lemmas}
\Crefname{originallemma}{Original lemma}{Original lemmas}
\makeatletter
\newlabel{perfect-lemma-resolution-property-perfect-complex}{{0F8F}{source}{}{sourcepointer-5}{}}
\newlabel{perfect-lemma-resolution-property-perfect-complex@cref}{{[originallemma][0][]0F8F}{[1][0][]source}{}{}{}}
\makeatother
\crefname{originalreference}{original reference}{original references}
\Crefname{originalreference}{Original reference}{Original references}
\makeatletter
\newlabel{situation-setup}{{02QL}{source}{}{sourcepointer-6}{}}
\newlabel{situation-setup@cref}{{[originalreference][0][]02QL}{[1][0][]source}{}{}{}}
\makeatother
\title{Sufficient conditions for the Chern classes of a perfect object to be defined}
\author{The Stacks Project Authors}
\date{stacks-project:chow.tex:local-snapshot-20261012; GNU FDL 1.2 or later}
\begin{document}
\maketitle
\noindent\textit{Editorial note.} All delivered mathematical text is the author's own, copied from the pinned Stacks Project chapter file chow.tex; no statement or proof step was rewritten, reordered or removed.\par\medskip
\noindent\textit{Editorial note.} Preamble provenance: the pinned chapter file opens with \textbackslash{}input\{preamble\} and carries no preamble text of its own, so the project's single Stacks-wide preamble staged at sources/stacks-chow/preamble.tex is excerpted verbatim as three spans, lines 8-12 (the verbatim package together with the project's own reference, slogan and history display-environment definitions, including the author's commented-out older form of the reference environment), lines 14-18 (Xy-pic setup) and lines 154-217 (the native theorem declarations and mathematical macros). The lines 8-12 span is required because a counted statement or its surrounding text may use those display environments, which the project typesets as comments. Nothing was added, removed or redefined; the wrapper class and the amsmath/amssymb/amsthm/graphicx base packages are supplied by the assembler.\par\medskip
\noindent\textit{Editorial note.} Same-chapter pointers are located at their \textbackslash{}label lines inside the pinned chapter file; cross-chapter pointers are located in the project's own tag list staged in the snapshot (sources/stacks-chow/tags.txt), at the verbatim lines carrying their canonical tags.\par\medskip
\noindent\textit{Editorial note.} The author's item (6) of the enumeration at source line 9316, 'add more here.', is retained verbatim as part of the statement; it is the author's own placeholder in the published text and was neither corrected nor removed. The phrase 'small detail omitted' at source line 9331 is likewise retained verbatim inside the counted proof.\par\medskip
\noindent\textit{Allowed standard prerequisites (editorial).}\begin{itemize}
\item The definition of the Chern classes of a perfect object. Same chapter chow.tex, label definition-defined-on-perfect part (1), invoked at proof line 9321. Identifies the first listed condition with the definition, so that it needs no argument.
\item Existence of an envelope with an ample family of invertible modules over a quasi-compact component. More on Morphisms, label more-morphisms-proposition-envelope-with-resolution-property, invoked at proof line 9330. Produces, for each irreducible component of X, an envelope of that component which has an ample family of invertible modules, which is the construction behind the implication from quasi-compact components to the first condition.
\item The resolution property and its consequence that a perfect complex can be represented by a bounded complex of finite locally free modules, plus base change of envelopes. Derived Categories of Schemes, labels perfect-lemma-resolution-property-ample-family (proof line 9333) and perfect-lemma-resolution-property-perfect-complex (proof line 9339); Same chapter chow.tex, label lemma-base-change-envelope (proof line 9349). The first two turn an ample family into the resolution property and then into a representing bounded complex, and the last shows that the base change of an envelope along the morphism of the fifth condition is again an envelope.
\end{itemize}\medskip
\ifdefined\hypertarget\hypertarget{sourcepointer-1}{}\fi
\noindent\textit{Source reference.} Original definition 0F9E: sources/stacks-chow/pinned-original-source.tex, source lines 9272--9272. Navigation pointer for \textbackslash{}ref\{definition-defined-on-perfect\} used inside the counted span (line 9321); the argument that consumes it is retained in the span.\par\medskip
\ifdefined\hypertarget\hypertarget{sourcepointer-2}{}\fi
\noindent\textit{Source reference.} Original lemma 0GU7: sources/stacks-chow/pinned-original-source.tex, source lines 3580--3580. Navigation pointer for \textbackslash{}ref\{lemma-base-change-envelope\} used inside the counted span (line 9350); the argument that consumes it is retained in the span.\par\medskip
\ifdefined\hypertarget\hypertarget{sourcepointer-3}{}\fi
\noindent\textit{Source reference.} Original proposition 0GTU: sources/stacks-chow/tags.txt, source lines 20979--20979. Navigation pointer for \textbackslash{}ref\{more-morphisms-proposition-envelope-with-resolution-property\} used inside the counted span (line 9330); the argument that consumes it is retained in the span.\par\medskip
\ifdefined\hypertarget\hypertarget{sourcepointer-4}{}\fi
\noindent\textit{Source reference.} Original lemma 0GMM: sources/stacks-chow/tags.txt, source lines 20762--20762. Navigation pointer for \textbackslash{}ref\{perfect-lemma-resolution-property-ample-family\} used inside the counted span (line 9333); the argument that consumes it is retained in the span.\par\medskip
\ifdefined\hypertarget\hypertarget{sourcepointer-5}{}\fi
\noindent\textit{Source reference.} Original lemma 0F8F: sources/stacks-chow/tags.txt, source lines 19040--19040. Navigation pointer for \textbackslash{}ref\{perfect-lemma-resolution-property-perfect-complex\} used inside the counted span (line 9339); the argument that consumes it is retained in the span.\par\medskip
\ifdefined\hypertarget\hypertarget{sourcepointer-6}{}\fi
\noindent\textit{Source reference.} Original reference 02QL: sources/stacks-chow/pinned-original-source.tex, source lines 1372--1372. Navigation pointer for \textbackslash{}ref\{situation-setup\} used inside the counted span (line 9301); the argument that consumes it is retained in the span.\par\medskip
\begin{lemma}
\label{lemma-chern-classes-defined}
Let $(S, \delta)$ be as in Situation \ref{situation-setup}.
Let $X$ be locally of finite type over $S$. Let $E \in D(\mathcal{O}_X)$
be a perfect object. If one of the following conditions hold, then
the Chern classes of $E$ are defined:
\begin{enumerate}
\item there exists an envelope $f : Y \to X$ such that $Lf^*E$
is isomorphic in $D(\mathcal{O}_Y)$ to a locally bounded complex of finite
locally free $\mathcal{O}_Y$-modules,
\item $E$ can be represented by a bounded complex of finite locally
free $\mathcal{O}_X$-modules,
\item the irreducible components of $X$ are quasi-compact,
\item $X$ is quasi-compact,
\item there exists a morphism $X \to X'$ of schemes locally of finite type
over $S$ such that $E$ is the pullback of a perfect object $E'$ on $X'$
whose chern classes are defined, or
\item add more here.
\end{enumerate}
\end{lemma}

\begin{proof}
Condition (1) is just Definition \ref{definition-defined-on-perfect} part (1).
Condition (2) implies (1).

\medskip\noindent
As in (3) assume the irreducible components $X_i$ of $X$ are quasi-compact.
We view $X_i$ as a reduced integral closed subscheme over $X$.
The morphism $\coprod X_i \to X$ is an envelope. For each $i$ there
exists an envelope $X'_i \to X_i$ such that $X'_i$ has an ample family
of invertible modules, see More on Morphisms, Proposition
\ref{more-morphisms-proposition-envelope-with-resolution-property}.
Observe that $f : Y = \coprod X'_i \to X$ is an envelope; small detail omitted.
By Derived Categories of Schemes, Lemma
\ref{perfect-lemma-resolution-property-ample-family}
each $X'_i$ has the resolution property.
Thus the perfect object $L(f|_{X'_i})^*E$ of $D(\mathcal{O}_{X'_i})$
can be represented by a bounded
complex of finite locally free $\mathcal{O}_{X'_i}$-modules, see
Derived Categories of Schemes, Lemma
\ref{perfect-lemma-resolution-property-perfect-complex}.
This proves (3) implies (1).

\medskip\noindent
Part (4) implies (3).

\medskip\noindent
Let $g : X \to X'$ and $E'$ be as in part (5). Then there exists an
envelope $f' : Y' \to X'$ such that $L(f')^*E'$ is represented by a
locally bounded complex $(\mathcal{E}')^\bullet$ of $\mathcal{O}_{Y'}$-modules.
Then the base change $f : Y \to X$ is an envelope by
Lemma \ref{lemma-base-change-envelope}. Moreover, the pulllback
$\mathcal{E}^\bullet = g^*(\mathcal{E}')^\bullet$ represents $Lf^*E$
and we see that the chern classes of $E$ are defined.
\end{proof}

\end{document}
